\begin{figure}[h]
\centering
\fbox{%
\begin{minipage}{0.94\linewidth}
\small
\textbf{\texttt{READER} prompt (abridged)}
\vspace{0.5em}

\ttfamily
You are reading a long document to answer a question.

You see one page this turn. Pages are ordered by retrieval relevance rather than document order. Previous pages are available only through NOTES.

\vspace{0.5em}
\textbf{WHAT TO RECORD}

Record what is on the CURRENT PAGE, and only evidence that bears on the QUESTION.

\begin{itemize}
    \item Preserve exact numbers, names, dates, identifiers, units, and relevant table structure.
    \item Describe relevant charts, figures, diagrams, and images in sufficient detail for later reasoning.
    \item Do not repeat evidence already present in NOTES.
    \item If the page adds nothing relevant, write \texttt{none}.
    \item If the QUESTION asks for a count or list, enumerate each matching item on this page and finish with:
\end{itemize}

\texttt{COUNT ON THIS PAGE: <number> (<what was counted>)}

Only items matching the QUESTION's exact definition should be counted.

\vspace{0.5em}
\textbf{OUTPUT}

\texttt{REASONING: <reasoning>}\\
\texttt{RECORD: <page evidence or none>}\\
\texttt{ANSWER PAGE: <true/false>}\\
\texttt{DECISION: <CONTINUE/STOP>}

\vspace{0.5em}
Write \texttt{STOP} only when the accumulated evidence appears sufficient to answer the QUESTION completely. Otherwise continue reading.
\end{minipage}%
}
\caption{Abridged prompt for the \texttt{reader} role. The production prompt additionally specifies formatting and output-recovery rules.}
\label{fig:method_reader_prompt}
\end{figure}